\scanpage{169}
\mspage{27}

on peut trouver

\[
z'\overset{\displaystyle\longrightarrow}{\deleted{\unsure{\longleftarrow|\longrightarrow}}}z
\]

égalisant la double-flèche $(*)$, donc \deleted{\unsure{OPS}} que celle-ci est déjà triviale. Appliquant encore une fois PS2, \deleted{\unsure{on trouve}} \inserted{par} un $z''\unsure{\mathbin{-}|\longrightarrow}z'$, on égalise de $=$ \unsure{la|le} double flèche $(*')$, et on gagne. On a donc prouvé, finalement

\diagram{\begin{tikzcd}[ampersand replacement=\&]
a \arrow[r,Leftrightarrow] \arrow[d,Rightarrow] \& a' \arrow[r,Rightarrow] \& b \arrow[r,Leftrightarrow] \& d\,(=\mathrm{PS1}+\mathrm{PS2})\\
c \arrow[r,Leftrightarrow] \& c' \& \&
\end{tikzcd}}

Mais pas l'implication inverse. En d'autres termes, si les topos induits $\unsure{\mathscr X|\mathcal X}_{/a}$, pour $a$ de $S$, sont \deleted{\unsure{on a expliqué}} tot. $0$-connexes, je \deleted{\unsure{ne}} n'ai pas prouvé qu'il en est de $=$ des $\mathscr X_{/W}$, $W$ un objet $0$-conn. \uline{quelc.} de $\mathscr X^\wedge$. \deleted{Notons que c) signifie \unsure{ici} que les produits fibrés de\unsure{ux}} \inserted{Ainsi, on \deleted{\unsure{est}} en est au point suivant}

\uline{\unsure{Proposition|Prop. Corollaire}} 2.11. Soit $X$ dans Cat. \deleted{Soit} Si $X$ est \inserted{loc} filtrante \inserted{\unsure{à g.|i.e.}}, alors \deleted{$X$ est $0$-connexe$,$} les produits fibrés $a\times_{\unsure{X|x}}b$ d'objets de $X$ sont des objets $0$-connexes de $X^\wedge$, \deleted{\unsure{ce qui}} \inserted{et cette} condition à son tour implique PS1, i.e. que $X$ est loc. irréductible. Mais en \unsure{partant|présence} de cette hypothèse PS1, ou $X$ loc irréductible, la condition sur les produits fibrés $a\times_{\unsure{x|\vphantom{x}}}b$ équivaut à \unsure{la|le} \unsure{suivante|suivant} : \deleted{pour tout \unsure{diagramme}} pour tout diagramme traits pleins

% archive sheet number: 169
